% Preámbulo común para las soluciones de los prácticos de Programación 3
\documentclass[11pt,a4paper]{article}

\usepackage[utf8]{inputenc}
\usepackage[T1]{fontenc}
\usepackage{lmodern}
\usepackage[spanish,es-noshorthands]{babel}
\usepackage[margin=2.3cm]{geometry}
\usepackage{amsmath,amssymb,amsthm}
\usepackage{enumitem}
\usepackage{xcolor}
\usepackage{tikz}
\usetikzlibrary{arrows.meta,positioning,shapes,calc}
\usepackage[noend]{algpseudocode}
\usepackage{algorithm}
\usepackage{booktabs}
\usepackage{tcolorbox}
\usepackage[hidelinks]{hyperref}

\floatname{algorithm}{Algoritmo}
\renewcommand{\algorithmicrequire}{\textbf{Entrada:}}
\renewcommand{\algorithmicensure}{\textbf{Salida:}}
\algrenewcommand\algorithmicwhile{\textbf{mientras}}
\algrenewcommand\algorithmicdo{\textbf{hacer}}
\algrenewcommand\algorithmicfor{\textbf{para}}
\algrenewcommand\algorithmicforall{\textbf{para todo}}
\algrenewcommand\algorithmicif{\textbf{si}}
\algrenewcommand\algorithmicthen{\textbf{entonces}}
\algrenewcommand\algorithmicelse{\textbf{sino}}
\algrenewcommand\algorithmicreturn{\textbf{retornar}}
\algrenewcommand\algorithmicfunction{\textbf{función}}
\algrenewcommand\algorithmicprocedure{\textbf{procedimiento}}

\theoremstyle{plain}
\newtheorem{lema}{Lema}[section]
\newtheorem{teorema}[lema]{Teorema}
\newtheorem{proposicion}[lema]{Proposición}
\theoremstyle{definition}
\newtheorem{definicion}[lema]{Definición}
\theoremstyle{remark}
\newtheorem*{observacion}{Observación}

% Caja para el enunciado de cada ejercicio
\newtcolorbox{enunciado}{colback=gray!6,colframe=gray!50,boxrule=0.4pt,arc=2pt,left=6pt,right=6pt,top=4pt,bottom=4pt,fontupper=\small}
% Nota que vincula el ejercicio con parciales/exámenes
\newcommand{\examref}[1]{\par\noindent{\small\color{blue!50!black}\textbf{Relación con parciales/exámenes:} #1}\par\medskip}

\newcommand{\NP}{\mathsf{NP}}
\newcommand{\Pclass}{\mathsf{P}}
\newcommand{\Oh}{\mathcal{O}}

\setlength{\parskip}{0.4em}
\setlength{\parindent}{0pt}


\title{Práctico 5 --- Algoritmos greedy}
\author{Programación 3 -- FING, UdelaR}
\date{2026}

\begin{document}
\maketitle

%=====================================================================
\section{Ejercicio 1}

\begin{enunciado}
\textbf{(a)} Dados un valor $N$, entero positivo, y una cantidad ilimitada de monedas de valores
10, 5, 2 y 1 pesos, presente un algoritmo que permita dar $N$ pesos usando una cantidad mínima
de monedas.
\textbf{(b)} Generalice el algoritmo para valores arbitrarios de las monedas.
\textbf{(c)} Si los valores de las monedas fueran 4 y 3, ¿el algoritmo seguiría dando la solución
óptima (esto es, con mínima cantidad de monedas)? ¿Y si los valores fueran 4, 3 y 1?
\end{enunciado}

\examref{Examen julio 2017, Pregunta~2 (b)--(c): el mismo algoritmo greedy de monedas, con
exactamente los contraejemplos de la parte (c) ($M=6$ con monedas $\{3,4\}$, donde no encuentra
solución, y con monedas $\{1,3,4\}$, donde encuentra una no óptima). También Examen febrero
2018, Pregunta~2 (contraejemplos con $\{3,5\}$, $N=9$ y con $\{1,5,8\}$, $N=10$).}

\subsection*{Parte (a)}
Una solución es una cuádrupla $(a,b,c,d)$ de naturales (cantidad de monedas de 10, 5, 2 y 1) con
$10a+5b+2c+d=N$; su costo es $a+b+c+d$. Siempre existe solución (por ejemplo $N$ monedas de 1) y
hay finitas, así que existe una óptima.

\paragraph{Algoritmo.} La elección greedy es: mientras falte dinero, dar la moneda de mayor valor
que no supere lo que falta. Como siempre se repite la misma moneda hasta que deja de caber, se
puede escribir con divisiones enteras:

\begin{algorithm}[H]
\caption{Cambio greedy con monedas de 10, 5, 2 y 1}
\begin{algorithmic}[1]
\Require $N\ge 1$
\Ensure $(a,b,c,d)$ cantidades de monedas de 10, 5, 2 y 1
\State $a\gets N \operatorname{div} 10$; \quad $N\gets N \bmod 10$
\State $b\gets N \operatorname{div} 5$; \quad $N\gets N \bmod 5$
\State $c\gets N \operatorname{div} 2$; \quad $N\gets N \bmod 2$
\State $d\gets N$
\State \Return $(a,b,c,d)$
\end{algorithmic}
\end{algorithm}
Como hay monedas de 1, el resto final siempre se completa y el resultado es una solución. El
tiempo de ejecución es $O(1)$ (cantidad constante de operaciones aritméticas); si se da moneda
por moneda, es $O(\text{cantidad de monedas})=O(N)$.

\paragraph{Correctitud.} Probamos que la solución óptima es \emph{única} y coincide con la
greedy.

\begin{lema}\label{lem:monedas}
Si $(a,b,c,d)$ es una solución óptima, entonces $d\le 1$, $c\le 2$, $b\le 1$ y no se cumple a
la vez $c=2$ y $d=1$. En consecuencia $2c+d\le 4$ y $5b+2c+d\le 9$.
\end{lema}
\begin{proof}
En cada caso, si no se cumpliera, podríamos reemplazar algunas monedas por otras de igual
suma total y menor cantidad, contradiciendo la optimalidad:
\begin{itemize}
  \item si $d\ge2$: cambiar dos monedas de 1 por una de 2 ($2$ monedas $\to 1$);
  \item si $c\ge3$: cambiar tres monedas de 2 por una de 5 y una de 1 ($3\to 2$);
  \item si $c\ge2$ y $d\ge1$: cambiar $2+2+1$ por una de 5 ($3\to1$);
  \item si $b\ge2$: cambiar dos monedas de 5 por una de 10 ($2\to1$).
\end{itemize}
Con $d\le1$, $c\le2$ y $(c,d)\ne(2,1)$, los valores posibles de $2c+d$ son $0,1,2,3,4$, así que
$2c+d\le4$; con $b\le1$ queda $5b+2c+d\le 9$.
\end{proof}

\begin{teorema}
El algoritmo devuelve la única solución óptima.
\end{teorema}
\begin{proof}
Sea $(a^*,b^*,c^*,d^*)$ una solución óptima. Como $N=10a^*+(5b^*+2c^*+d^*)$ con
$0\le 5b^*+2c^*+d^*\le 9$ (lema), por la unicidad del cociente y el resto de la división entera,
$a^*=N\operatorname{div}10$ y $5b^*+2c^*+d^*=N\bmod 10$. Llamemos $r_1=N\bmod 10$. Como
$r_1=5b^*+(2c^*+d^*)$ con $0\le 2c^*+d^*\le4$, se tiene $b^*=r_1\operatorname{div}5$ y
$2c^*+d^*=r_1\bmod 5=:r_2$. Finalmente $r_2=2c^*+d^*$ con $0\le d^*\le1$, así que
$c^*=r_2\operatorname{div}2$ y $d^*=r_2\bmod 2$. Estos son exactamente los valores que calcula
el algoritmo en los pasos 1--4. Luego la solución greedy es igual a toda solución óptima.
\end{proof}

\subsection*{Parte (b)}
Sean $v_1,\dots,v_k$ los valores (enteros positivos y distintos) de las monedas. El algoritmo
generalizado ordena los valores de mayor a menor y aplica la misma elección greedy: usar tantas
monedas como entren de la mayor denominación que no supere el monto restante.

\begin{algorithm}[H]
\caption{Cambio greedy con valores arbitrarios}
\begin{algorithmic}[1]
\Require $N\ge1$; valores $v_1,\dots,v_k$
\Ensure cantidades $(c_1,\dots,c_k)$, o ``no encontró solución''
\State Ordenar los valores de modo que $v_1>v_2>\cdots>v_k$
\For{$i$ desde $1$ hasta $k$}
  \State $c_i\gets N\operatorname{div} v_i$
  \State $N\gets N\bmod v_i$
\EndFor
\If{$N=0$} \State \Return $(c_1,\dots,c_k)$
\Else \State \Return ``no encontró solución''
\EndIf
\end{algorithmic}
\end{algorithm}

El tiempo de ejecución es $O(k\log k)$ por el ordenamiento más $O(k)$ por el ciclo. Si alguna
moneda vale 1, el algoritmo siempre devuelve una solución ($N\bmod 1=0$). Sin embargo, como
muestra la parte (c), para valores arbitrarios el algoritmo \textbf{no} garantiza encontrar una
solución ni que la que encuentre sea óptima; la demostración de la parte (a) usa propiedades
particulares del sistema $\{10,5,2,1\}$ (el Lema~\ref{lem:monedas}).

\begin{observacion}
Para valores arbitrarios el problema se resuelve correctamente con programación dinámica:
$\mathrm{OPT}(0)=0$ y $\mathrm{OPT}(x)=1+\min\{\mathrm{OPT}(x-v_i): v_i\le x\}$ (con
$\min\emptyset=\infty$), en tiempo $O(Nk)$, que es pseudo-polinomial.
\end{observacion}

\subsection*{Parte (c)}
\textbf{Monedas de 4 y 3: no.} Con $N=6$, el algoritmo toma una moneda de 4 ($6\operatorname{div}4=1$)
y queda $N=2$; como $2<3$ no puede dar más monedas y termina con resto $2\ne 0$: no encuentra
solución. Pero existe la solución óptima $3+3$ (dos monedas). (Otro ejemplo: $N=9$, el algoritmo
da $4+4$ y queda resto 1, mientras que $3+3+3$ es solución.)

\textbf{Monedas de 4, 3 y 1: tampoco.} Ahora siempre encuentra una solución (hay moneda de 1),
pero no necesariamente óptima. Con $N=6$ el algoritmo da $4+1+1$ (tres monedas), mientras que
$3+3$ usa dos monedas.

%=====================================================================
\section{Ejercicio 2}

\begin{enunciado}
(K\&T, Ex.~4.5) Hay casas a lo largo de una carretera recta, cada una dada por su distancia
$x_i$ al inicio de la carretera. Se quieren ubicar estaciones de modo que ninguna casa esté a más
de 4 km de una estación.
\textbf{(a)} Dé un algoritmo eficiente que determine la ubicación de las estaciones de manera que
se satisfaga el requerimiento de cobertura con la menor cantidad posible de estaciones.
\textbf{(b)} Demuestre la corrección de su algoritmo. \emph{Sugerencia:} siga una estrategia
similar a la de la sección 4.1 del libro.
\end{enunciado}

\examref{Es el mismo problema que el Primer parcial 2020, Ej.~2 (tiempos de emisión de un
mensaje que cubran todas las clases) y el Primer parcial 2023, Ej.~3 (horas de entrega del
carpintero): cada casa $x$ define el intervalo $[x-4,x+4]$ de posiciones de estación que la
cubren, y se busca un conjunto mínimo de puntos que \emph{pinche} todos los intervalos. Nuestro
algoritmo es el de las soluciones oficiales (elegir el menor extremo derecho entre los intervalos
no cubiertos) y la demostración sigue la misma estructura ``greedy está por delante''.}

\subsection*{Parte (a)}
Decimos que una estación en la posición $p$ \emph{cubre} la casa $x$ si $|x-p|\le 4$. Idea
greedy: la casa más a la izquierda que todavía no está cubierta tiene que ser cubierta por alguna
estación; la ponemos \emph{lo más a la derecha posible}, a 4 km a su derecha, para cubrir la
mayor cantidad posible de casas siguientes.

\begin{algorithm}[H]
\caption{Ubicación greedy de estaciones}
\begin{algorithmic}[1]
\Require posiciones de las casas $x_1,\dots,x_n$
\Ensure posiciones de las estaciones $s_1<\dots<s_k$
\State Ordenar las casas de modo que $x_1\le x_2\le\cdots\le x_n$
\State $i\gets 1$; \ $k\gets 0$
\While{$i\le n$} \Comment{$x_i$ es la casa no cubierta más a la izquierda}
  \State $k\gets k+1$; \ $s_k\gets x_i+4$
  \While{$i\le n$ \textbf{y} $x_i\le s_k+4$} \Comment{marcar como cubiertas}
    \State $i\gets i+1$
  \EndWhile
\EndWhile
\State \Return $s_1,\dots,s_k$
\end{algorithmic}
\end{algorithm}

\paragraph{Tiempo de ejecución.} El ordenamiento cuesta $O(n\log n)$. En el resto, el índice
$i$ sólo aumenta y lo hace a lo sumo $n$ veces; cada iteración del ciclo externo incrementa $i$
al menos una vez (la casa $x_i$ con la que se define $s_k$ cumple $x_i\le s_k+4$). Por lo tanto
los ciclos cuestan $O(n)$ y el total es $O(n\log n)$ ($O(n)$ si las casas vienen ordenadas).

\subsection*{Parte (b)}
Para cada $j$, sea $y_j$ la casa $x_i$ con la que se define $s_j=y_j+4$ en el paso 4.

\begin{lema}\label{lem:valida}
\begin{enumerate}[label=(\alph*)]
  \item Las casas que el ciclo interno marca en la iteración $j$ están cubiertas por $s_j$.
  \item Al comenzar la iteración $j+1$ se cumple $y_{j+1}>s_j+4$. En particular
        $s_1<s_2<\dots<s_k$.
  \item El algoritmo termina y toda casa queda cubierta: el resultado es una solución.
\end{enumerate}
\end{lema}
\begin{proof}
(a) Las casas marcadas en la iteración $j$ son las $x$ con $y_j\le x\le s_j+4$ (las casas están
ordenadas y se marcan a partir de $y_j$). Como $s_j=y_j+4$, se tiene $s_j-4\le x\le s_j+4$, es
decir $|x-s_j|\le 4$.
(b) El ciclo interno de la iteración $j$ termina con $i>n$ o con $x_i>s_j+4$; si comienza la
iteración $j+1$ es porque $i\le n$, así que $y_{j+1}=x_i>s_j+4$. Luego
$s_{j+1}=y_{j+1}+4>s_j+8>s_j$.
(c) Cada iteración del ciclo externo incrementa $i$ (ver el análisis de tiempo), así que el
algoritmo termina, y lo hace con $i>n$, es decir, con todas las casas marcadas; por (a) todas
están cubiertas.
\end{proof}

Sea ahora $O=\{o_1<o_2<\dots<o_m\}$ una solución cualquiera (posiciones distintas, ordenadas de
menor a mayor; si una solución repitiera posiciones, eliminando repeticiones se obtiene otra con
menos estaciones). Mostramos que el algoritmo está ``por delante'' de $O$: sus estaciones están
siempre a la derecha (o en el mismo lugar) que las de $O$, de modo que cubren al menos lo mismo.

\begin{proposicion}\label{prop:adelante}
Para todo $j$, $1\le j\le\min\{k,m\}$, se cumple $o_j\le s_j$.
\end{proposicion}
\begin{proof}
Por inducción en $j$.

\emph{Caso $j=1$.} Tenemos $y_1=x_1$ y $s_1=x_1+4$. Como $O$ es solución, la casa $x_1$ está
cubierta por alguna $o_l$, de modo que $o_l\le x_1+4$. Como $o_1\le o_l$, resulta
$o_1\le x_1+4=s_1$.

\emph{Paso inductivo.} Supongamos $o_j\le s_j$ para cierto $j<\min\{k,m\}$. Por el
Lema~\ref{lem:valida}(b) y la hipótesis de inducción,
\[
y_{j+1}>s_j+4\ge o_j+4\ge o_l+4\qquad\text{para todo } l\le j,
\]
así que ninguna de las estaciones $o_1,\dots,o_j$ cubre la casa $y_{j+1}$. Como $O$ es
solución, $y_{j+1}$ está cubierta por alguna $o_l$ con $l\ge j+1$, y entonces
$o_{j+1}\le o_l\le y_{j+1}+4=s_{j+1}$.
\end{proof}

\begin{teorema}
El algoritmo devuelve una solución con la mínima cantidad de estaciones.
\end{teorema}
\begin{proof}
Por el Lema~\ref{lem:valida}(c) el resultado es una solución. Sea $O$ una solución
cualquiera con $m$ estaciones, como antes; veamos que $k\le m$. Supongamos por absurdo que
$k>m$. Entonces existe la iteración $m+1$ y, por el Lema~\ref{lem:valida}(b) y la
Proposición~\ref{prop:adelante} con $j=m$,
\[
y_{m+1}>s_m+4\ge o_m+4\ge o_l+4\qquad\text{para todo } l\le m,
\]
por lo que la casa $y_{m+1}$ no está cubierta por ninguna estación de $O$, lo que contradice que
$O$ sea solución. Luego $k\le m$ y, como $O$ es arbitraria, la solución greedy es óptima.
\end{proof}

\begin{observacion}
Si la carretera termina en la posición $L$ (con todas las casas en $[0,L]$) y las estaciones deben
estar sobre la carretera, basta usar $s_k\gets\min\{x_i+4,\,L\}$ en el paso 4. La demostración es
la misma: si $s_k=L$ todas las casas restantes están en $[L-4,L]$ y quedan cubiertas, y en la
proposición se usa además que $o_l\le L$.
\end{observacion}

%=====================================================================
\section{Ejercicio 3}

\begin{enunciado}
(K\&T, Ex.~4.6) En un triatlón, para cada participante $i$ se tiene una estimación del tiempo de
natación $s_i$, de ciclismo $b_i$ y de carrera $r_i$. La piscina sólo puede ser usada por un
nadador a la vez; las etapas posteriores pueden desarrollarse simultáneamente por varios
competidores. Se quiere determinar en qué orden largan los competidores (uno a la vez) para
minimizar la duración de la competencia, que termina cuando el último participante llega a la
meta.
\textbf{(a)} Dé un algoritmo eficiente para resolver el problema.
\textbf{(b)} Demuestre la corrección de su algoritmo.
\end{enunciado}

\examref{El Primer parcial 2022, Ej.~2 (panadería: preparación en paralelo y horneado de a uno)
es la versión ``espejada'' de este problema: allí la etapa compartida es la última y se ordena
por preparación creciente; aquí la etapa compartida es la primera y se ordena por
$b_i+r_i$ decreciente. Seguimos la misma técnica de la solución oficial: \emph{exchange
argument} sobre inversiones entre posiciones contiguas.}

\paragraph{Modelo.} Sea $t_i=b_i+r_i$ el tiempo que le lleva al participante $i$ terminar
una vez que sale de la piscina. Una solución es un orden $O=(o_1,\dots,o_n)$ de largada. Podemos
suponer que la piscina nunca queda ociosa: cada nadador entra apenas sale el anterior (dejar la
piscina vacía un tiempo sólo atrasa a todos los que vienen después y no adelanta a nadie). Así,
el participante $o_j$ sale de la piscina en el instante $\sum_{l\le j}s_{o_l}$ y llega a la
meta en
\[
F_j(O)=\sum_{l=1}^{j} s_{o_l}+t_{o_j},
\qquad\text{y la duración de la competencia es}\qquad
T(O)=\max_{1\le j\le n} F_j(O).
\]
Queremos un orden $O$ que minimice $T(O)$.

\subsection*{Parte (a)}
Hacer largar primero a los que más tardan en las etapas que se hacen en paralelo.

\begin{algorithm}[H]
\caption{Orden de largada del triatlón}
\begin{algorithmic}[1]
\Require $(s_i,b_i,r_i)$ para $i=1,\dots,n$
\Ensure orden de largada $G=(g_1,\dots,g_n)$
\State Calcular $t_i=b_i+r_i$ para cada $i$
\State Ordenar los participantes de modo que $t_{g_1}\ge t_{g_2}\ge\cdots\ge t_{g_n}$ (empates en cualquier orden)
\State \Return $G$
\end{algorithmic}
\end{algorithm}
El tiempo de ejecución es $O(n)$ para el paso 1 más $O(n\log n)$ para el ordenamiento:
$O(n\log n)$.

\paragraph{Otros criterios naturales no funcionan.} Por ejemplo, ordenar por natación
creciente: con $1$: $(s,t)=(1,1)$ y $2$: $(s,t)=(10,100)$, el orden $(1,2)$ da
$T=\max\{2,\,111\}=111$, mientras que $(2,1)$ da $T=\max\{110,\,12\}=110$. Ordenar sólo por
carrera $r_i$ decreciente: con $1$: $(s,b,r)=(1,10,0)$ y $2$: $(s,b,r)=(1,0,1)$, ese criterio
da $(2,1)$ con $T=\max\{2,12\}=12$, pero $(1,2)$ da $T=\max\{11,3\}=11$.

\subsection*{Parte (b)}
Fijemos el orden $G=(g_1,\dots,g_n)$ que devuelve el algoritmo. Dado otro orden $O$, decimos que
dos participantes $p,q$ forman una \emph{inversión} de $O$ (respecto de $G$) si aparecen en
distinto orden relativo en $O$ y en $G$. $G$ es el único orden sin inversiones.

\begin{lema}\label{lem:contigua}
Si $O\neq G$, entonces existen dos participantes \emph{contiguos} en $O$, $o_j$ y $o_{j+1}$, que
forman una inversión (es decir, $o_{j+1}$ está antes que $o_j$ en $G$).
\end{lema}
\begin{proof}
Sea $\pi(p)$ la posición de $p$ en $G$. Si no hubiera inversiones contiguas, se cumpliría
$\pi(o_1)<\pi(o_2)<\dots<\pi(o_n)$; como $\pi$ es una biyección sobre $\{1,\dots,n\}$, esto
implica $\pi(o_j)=j$ para todo $j$, es decir $O=G$.
\end{proof}

\begin{lema}\label{lem:intercambio}
Sean $p=o_j$ y $q=o_{j+1}$ contiguos en $O$ formando una inversión, y sea $O'$ el orden que
resulta de intercambiarlos. Entonces $T(O')\le T(O)$, y $O'$ tiene exactamente una inversión
menos que $O$.
\end{lema}
\begin{proof}
Como $q$ está antes que $p$ en $G$ y $G$ está ordenado por $t$ decreciente, $t_q\ge t_p$. Sea
$A=\sum_{l<j}s_{o_l}$ el tiempo de natación de los que largan antes que ambos.
Para los participantes en posiciones $l<j$ o $l>j+1$ la suma de tiempos de natación hasta ellos
inclusive no cambia (es la misma suma de términos), así que su tiempo de llegada es el mismo en
$O$ y en $O'$. Para $p$ y $q$:
\[
\begin{array}{ll}
\text{en } O: & F_p=A+s_p+t_p,\qquad F_q=A+s_p+s_q+t_q,\\[2pt]
\text{en } O': & F'_q=A+s_q+t_q,\qquad F'_p=A+s_q+s_p+t_p .
\end{array}
\]
Entonces $F'_q\le F_q$ (porque $s_p\ge0$) y $F'_p=A+s_p+s_q+t_p\le A+s_p+s_q+t_q=F_q$ (porque
$t_p\le t_q$). Así $\max\{F'_p,F'_q\}\le F_q\le T(O)$ y, como los demás tiempos de llegada no
cambian, $T(O')\le T(O)$.

Por último, el intercambio de dos elementos contiguos sólo cambia el orden relativo del par
$\{p,q\}$, que pasa a coincidir con el de $G$; el orden relativo de cualquier otro par no cambia.
Luego $O'$ tiene una inversión menos.
\end{proof}

\begin{teorema}
El orden $G$ devuelto por el algoritmo es óptimo.
\end{teorema}
\begin{proof}
Sea $O^*$ un orden óptimo (existe porque hay finitos órdenes). Si $O^*=G$ no hay nada que
probar. Si no, por el Lema~\ref{lem:contigua} tiene una inversión contigua y por el
Lema~\ref{lem:intercambio} al intercambiarla obtenemos un orden con una inversión menos y
duración no mayor. Repitiendo, como la cantidad de inversiones es un natural (a lo sumo
$\binom n2$) que disminuye en uno en cada paso, después de finitos intercambios llegamos a un
orden sin inversiones, que es $G$, y en cada paso la duración no aumentó:
\[
T(G)\le\cdots\le T(O^*).
\]
Como $O^*$ es óptimo, $T(G)=T(O^*)$ y $G$ también es óptimo.
\end{proof}

\end{document}
